\chapter{The Imaginary Bounce}

This is the chapter where the instrument sticks its neck out. Everything the book has sealed --- the sideways
direction the corridor carried and could not cash, the phase face that was decided but not explained, the sign the
naming read as a minus and would not interpret --- comes out here, and comes out as a single bold claim about the
world. But the boldness must be handled with the most exacting honesty in the volume, because this chapter is
where it would be easiest, and most fatal, to let a beautiful reading wear the certainty of a plain proof. So the
chapter does two things at once and keeps them ruthlessly apart: it makes the prediction, in full, with all its
imaginative reach; and it grades the prediction, line by line, so that you always know exactly which part is
proved, which part is a reading offered over the proof, and which part is reserved entirely for a hand the
construction does not contain. Read the reveal for its daring. Read the grading for its honesty. The book stakes
its whole method on your being able to tell the two apart.

\section{The Direction Cashed}

Begin with what the corridor hid. Chapters ago the machine carried a direction it could not read as a value --- a
something-to-the-side of everything countable, present and uncashable, given a face on the number as the phase but
never explained. Here is the reading the book has been saving for it: that sideways direction is the
\emph{imaginary} axis. It runs off the real, countable line at a right angle, and the book reads it as the machine
carrying, alongside its real magnitudes, a component in imaginary time --- the axis a physicist reaches for when a
rotation has to be represented as something other than a stretch. The uncashable direction of the corridor is
cashed at last, and the coin it is cashed for is the imaginary unit.

Say clearly, at the very moment of the reveal, what kind of claim this is, because everything downstream depends
on it. That the machine carries a sideways direction is built --- the corridor really does hold an orientation off
the countable line, and that is construction, not interpretation. That this sideways direction \emph{is} the
imaginary axis, imaginary time, the $-i$ a rotation turns through --- that is a reading. It is the book's
interpretation of what the built sideways direction amounts to when set against physics, offered as the most
natural reading of a real structural feature, but not a thing the construction proves. The direction is built; its
identification as the imaginary axis is a reading laid over the building. Hold that seam; the whole chapter is
about keeping it visible.

\section{The Bounce, and the Square}

Now the physics the reading opens, told as a reading throughout. Picture two measurements set at right angles ---
crossed at ninety degrees, the way a polarizer crossed against another is supposed to extinguish the light, or the
way two settings of a quantum measurement can be made maximally incompatible. The naive expectation is that at a
right angle the amplitude simply dies: crossed filters go dark, the reading falls to nothing. The book's reading
says otherwise. At the crossed right angle the amplitude does not extinguish; it \emph{rotates}. It turns through
the sideways direction just cashed --- multiplied by the imaginary unit, sent a quarter-turn into imaginary time
rather than driven to zero. The bounce, the book calls it: at ninety degrees the reading does not vanish, it
bounces sideways, $\times(-i)$, the turn a physicist knows as the rotation into imaginary time.

The picture is worth drawing concretely, because it is where the reading touches physics a reader may already
know. Two polarizing filters crossed at a right angle pass no light --- that is the textbook fact, the darkness at
ninety degrees. But the quantum fills the crossed angle with something the classical picture misses: slip a third
filter between the two, at forty-five degrees, and light reappears beyond the pair that alone let nothing through.
The between-angle added no light; it \emph{rotated} the amplitude, turned it partway rather than extinguishing it,
so that what the crossed pair killed the intermediate turn revived. The book reads that reappearance as the
bounce: the amplitude at an angle does not die, it turns, and the turning is through the sideways direction the
book has cashed as imaginary. The same shape shows in the wave a moving particle carries --- the phase that
advances not as a growing magnitude but as a rotation, a hand sweeping round rather than a line getting longer.
Rotation, not extinction; a turn through the imaginary, not a fall to zero. That is the physical face of the
reading, offered as the natural interpretation of a proved two-fold and, still, a reading.

And here the sign the naming would not interpret finds its meaning, as a reading. Take that imaginary quarter-turn
and do it twice --- bounce, and bounce again --- and you have turned through the imaginary direction a full half
circle, and a half circle is a reversal: $(-i)$ squared is $-1$. So the book reads the proved minus of the last
chapter --- the sign the electron's charge was read to carry --- as exactly this: the square of an imaginary
quarter-turn, $-1 = (-i)^2$, two bounces through imaginary time. The one side of the split is the thing turned once
too few or once too many; the reversal between them, the flip from plus to minus, is read as two imaginary
bounces. And carry the turning further --- four quarter-turns, $\{1, -i, -1, i\}$ --- and you come back to where you
started only after going around not once but, in the doubled bookkeeping the imaginary axis forces, the full
figure a physicist reads as a particle that must turn through twice as much as a plain arrow to come home. The
four-fold of the imaginary unit is read as that doubled return.

Every sentence of the last two paragraphs is a reading, and the book will not let one of them pass as a proof.
There is no theorem here that the amplitude rotates by the imaginary unit, no demonstration that the sign is a
squared quarter-turn, no proof that the four-fold is the doubled return. These are the book's interpretation of a
proved structure --- the interpretive bridge from an integer the machine genuinely established to the imaginary,
rotational, physical picture the integer invites. The bridge is beautiful and, I will argue, apt; it is not
proved, and the next section says exactly what \emph{is}.

\section{The Proved Flip, and the Read Bridge}

Here is the grade, split as finely as the volume knows how, because this split is the whole discipline of the book
brought to its sharpest point.

What is PROVED is the flip. Drive the loop over the symmetric baseline and over a tilted one, and the sign the
recovered number carries flips --- one way on the flat, the other way on the tilt, a clean two-fold reversal. That
flip is a theorem, and it is the cleanest theorem in the volume: it depends on \emph{no axioms at all}, decided by
finite counting alone, forced from the one proved floor the pigeonhole named the electron in. The integer
two-fold --- that there are two signs, that they are genuinely opposite, that the construction flips between them
--- is proved, with the strongest grade the book can give, nothing assumed, nothing granted, pure counting. Say it
without hedging: the $\pm 1$ flip is proved.
% CHOICE (vol1-interp): the +-1 flip is PROVED (baseline_relative_flip = [], no axioms). The identification of that
% flip with (-i)^2, of the sideways direction with the -i axis / imaginary time, of the four-fold with the spinor's
% 4pi return, and of the two signs with matter/antimatter, is the INTERPRETIVE BRIDGE -- a reading laid over the
% proved flip, offered and apt but never proved. Do not let the -i reading wear the flip's PROVED.

What is a READING --- offered, marked, never proved --- is everything the last two sections built on top of that
flip. That the two signs are $+1$ and $(-i)^2$; that the sideways direction is the imaginary axis; that the
crossed-angle amplitude bounces by the imaginary unit; that the four-fold is a doubled return; that the whole
two-fold is to be read as the split between matter and its opposite --- all of that is the interpretive bridge,
the book's reading of what the proved flip means when carried into physics. It is the boldest reading in the
volume and it is still a reading. The construction proves that there are two opposite signs; the book reads those
two signs as matter and antimatter, and reads the reversal between them as a squared imaginary quarter-turn. The
proof gives you the integer; the reading gives you the picture; and the book insists, here at its most tempting
moment, that the picture does not inherit the integer's certainty. The flip is proved. The bridge is read. Keep
them apart and the chapter is honest; blur them and it is a conjuring trick with the volume's best trick.

\section{The Prediction, the Constant, and the Reader}

With the grade kept straight, the book can now state its prediction plainly, because a prediction stated with its
grade attached is a claim you can trust rather than a flourish you must swallow. Vol 1 predicts the split: that the
world's smallest matter comes with a two-fold sign, a proved flip read as matter against antimatter, the reversal
between them read as two turns through imaginary time. The integer skeleton of that prediction is proved; the
imaginary flesh on it is the book's reading; and together they are the neck the volume sticks out --- a claim
about the world derived from a finite construction, offered for the world to confirm or break.

It is worth being clear about what kind of prediction this is, because a prediction graded this carefully is a
different and stronger thing than a bold guess. The book does not predict the split the way one predicts a horse
race, by intuition dressed as inference. It predicts the \emph{shape} of the split by derivation --- the proved
integer two-fold --- and then reads that shape into the physical picture of matter and its mirror. A reader who
accepts only the proof is left with a genuine, unhedged result: the world's smallest matter carries a two-fold
sign, established by counting, owing nothing to interpretation. A reader who accepts the reading as well is offered
the fuller picture: that two-fold as the matter/antimatter split, the reversal as a squared turn through imaginary
time. The book gives both and marks the seam between them, so that neither reader is misled --- the cautious one
gets a proof with no reading smuggled in, the venturesome one a reading that never pretends to be the proof. That
is what it is to stick a neck out honestly: to make the boldest claim the construction supports, and to label, to
the word, exactly how much of it is earned.

One step of the prediction the construction deliberately does not take, and the reservation is as important as the
reveal. A physical split in the world comes with a \emph{number} --- a dimensionful constant, a measured size, the
actual magnitude at which the world sets the thing the book has predicted the shape of. That last step, from the
proved-and-read \emph{shape} of the split to an actual physical constant with units on it, the construction does
not make. It builds all the way up to the split --- the two-fold, the sign, the read as matter and antimatter ---
and it stops there, at the shape, and hands the closing of the shape to a number over to a determination it does
not itself contain. The book will not invent that number. It stakes the split and reserves the constant, and it
says so in the open.
% CHOICE (operator): the dimensionful physical constant that closes the predicted split to an actual measured
% number is RESERVED -- the construction builds up to the split (+-1 = (-i)^2, matter/antimatter shape) and does
% NOT invent the constant; the last step to a number is the operator's.

And one thing more the construction hands over, not to an operator but to you. Which side of the proved flip is
``matter'' and which is ``antimatter'' is not fixed by the construction. The two signs are genuinely opposite ---
that is proved --- but which of the opposites wears which name is a matter of orientation, of the frame the reader
stands in to look at the split, and the construction is scrupulously silent on it. It gives you the two-fold and
the flip and the reading; it does not give you which way round to read the sign, because that is the reader's to
set. So the volume's boldest chapter ends, characteristically, by handing something back: the split proved in its
integer, read in its imaginary flesh, reserved in its constant, and left --- as to which side is which --- to the
one holding the book. The neck is stuck out. What the world does with it, and which way you choose to face it, the
construction leaves open on purpose.
